% SPDX-License-Identifier: Apache-2.0
\section{Four Outputs Between On and Off}
\label{sec:x-pwm}

A pin is on or off, and much of what a board drives wants a value in
between: an LED at a third of its brightness, a fan at half its speed,
a servo told where to stand. Switching the pin fast enough that the
load cannot follow gives the in-between, and what the load sees is the
fraction of each period the pin was high.

\code{Pwm} is four such outputs behind AXI-Lite, in six words: a
period, a duty per channel, and a control word with an enable, a
centre bit and a polarity bit per channel. A counter runs the period,
and each channel is high while the counter is below its duty.

The ends need no special case, which is the reason for that
comparison. A duty of zero never wins it, so the pin never rises; a
duty of the period or more always wins it, so the pin never falls.
Neither end is a pulse one cycle wide, which is what a scheme that
counted edges rather than compared would leave behind.

A duty written takes effect at the end of the period rather than at
once. A width that changed under the counter halfway through would
make one pulse of neither the old width nor the new, and a load that
averages sees a step nobody asked for while a servo reads it as a
command. So a write goes to a shadow, and every shadow is loaded
together at the wrap. The run below widens a duty from two cycles to
eight in the middle of a period, and every whole pulse it then
measures is two wide or eight, never something between them.

Two shapes of pulse. Edge aligned counts up and wraps, so every pulse
starts at the same moment and its end moves as the duty does. Centre
aligned counts up and then back down, so the pulse is centred in its
period and both edges move; a bank of outputs driven that way does not
switch all at once, which is what a supply prefers. A centred period
is twice the count, and the run shows channel 2, at four of ten, as
one run of eight cycles rather than two runs at the ends of the
period.

\begin{figure}[htbp]
\centering
\input{pwm_timing.tex}
\caption{The counter and the four pins, through one change of duty:
\code{count} runs the period and each pin is high below its duty.}
\label{fig:x-pwm}
\end{figure}

\lstinputlisting[caption={\code{ex\_pwm.rs}. Run with \code{bazel run //lib/examples:ex\_pwm}.},label={lst:x-pwm}]{lib/examples/ex_pwm.rs}

\lstinputlisting[style=ascii,caption={What \code{ex\_pwm} printed when the build ran it: the duties as cycles high, the pulses around a change of width, the channel turned over, and the centred pulses.},label={lst:x-pwm-out}]{out_pwm.txt}
